\chapter{MÔ HÌNH HỆ THỐNG VÀ BÀI TOÁN TỐI ƯU}
\label{chap:mo-hinh-he-thong}

\section{Kiến trúc hệ thống}
\label{sec:kien-truc-he-thong}

Luận văn xét một hệ thống truyền xuống gồm một BS SISO, một STAR--RIS thụ động có $N$
phần tử và $K = 4$ UE SISO. Các UE được bố trí ở hai phía của STAR--RIS. Trong mã nguồn,
chỉ báo phía người dùng được gán luân phiên, vì vậy hai UE thuộc phía phản xạ và hai UE
thuộc phía truyền. Mỗi thời điểm, BS phát đồng thời một luồng chung và $K$ luồng riêng
theo kiến trúc RSMA một lớp.

Hình~\ref{fig:ch2-mo-hinh} minh họa các đường truyền chính. Một UE nhận tín hiệu từ hai
nguồn: đường trực tiếp BS--UE và đường ghép tầng BS--STAR--RIS--UE. STAR--RIS không tự
phát dữ liệu mà điều khiển biên độ tương đối và pha của tín hiệu đi qua hoặc phản xạ tại
bề mặt. Vì vậy, hiệu quả của STAR--RIS thể hiện qua việc làm thay đổi kênh hiệu dụng mà
RSMA nhìn thấy. Trong hình, đường liền nét là đường ghép tầng qua bề mặt, còn đường nét
đứt là đường trực tiếp; đường trực tiếp tồn tại với mọi UE, kể cả các UE nằm ở phía
truyền qua.

\begin{figure}[htbp]
\centering
\resizebox{0.95\textwidth}{!}{\begin{tikzpicture}[font=\scriptsize,>=Stealth,line width=0.5pt]

% --- BS mot anten ---
\node[bx,minimum height=0.85cm,minimum width=0.55cm] (bs) at (0.45,1.85) {BS};
\draw[cUEc,line width=0.6pt] (0.72,1.85) -- ++(0.20,0);
\fill[cUEc] (0.92,1.78) -- (0.92,1.92) -- (1.08,1.85) -- cycle;

% --- STAR-RIS ---
\foreach \i in {0,...,10}{
   \fill[cKenh!22] (4.29,0.42+\i*0.28) rectangle ++(0.22,0.22);
   \draw[cKenh,line width=0.35pt] (4.29,0.42+\i*0.28) rectangle ++(0.22,0.22);
}
\draw[cKenh,line width=0.8pt] (4.40,0.34) -- (4.40,3.62);
\node[font=\tiny,anchor=north,align=center] at (4.40,0.26) {STAR--RIS\\ $N$ phần tử};

% --- UE ---
\node[ue,fill=cUEa!18,draw=cUEa] (u1) at (2.25,3.45) {UE$_1$};
\node[ue,fill=cUEc!18,draw=cUEc] (u3) at (2.45,0.55) {UE$_3$};
\node[ue,fill=cUEb!18,draw=cUEb] (u2) at (6.75,3.45) {UE$_2$};
\node[ue,fill=cUEd!18,draw=cUEd] (u4) at (6.95,0.55) {UE$_4$};

% --- duong BS -> STAR-RIS : g ---
\draw[cKenh,line width=0.9pt,->] (1.10,1.85) -- (4.25,2.00);
\node[font=\scriptsize,cKenh,anchor=south] at (2.75,2.02) {$g$};

% --- duong ghep tang STAR-RIS -> UE : h_{r,k} ---
\draw[cUEa,->] (4.32,2.55) -- (2.55,3.28);
\node[font=\tiny,cUEa,anchor=south west] at (3.35,3.00) {$h_{r,1}$};
\draw[cUEc,->] (4.32,1.35) -- (2.72,0.72);
\node[font=\tiny,cUEc,anchor=north west] at (3.35,1.02) {$h_{r,3}$};
\draw[cUEb,->] (4.52,2.55) -- (6.45,3.28);
\node[font=\tiny,cUEb,anchor=south east] at (5.85,3.00) {$h_{r,2}$};
\draw[cUEd,->] (4.52,1.35) -- (6.65,0.72);
\node[font=\tiny,cUEd,anchor=north east] at (5.85,1.02) {$h_{r,4}$};

% --- duong truc tiep BS -> UE : h_{d,k} ---
\draw[gray!60,dashed,->] (1.05,2.10) -- (2.05,3.22);
\node[font=\tiny,gray!45!black,anchor=south] at (1.52,2.98) {$h_{d,1}$};
\draw[gray!60,dashed,->] (1.05,1.45) -- (2.25,0.78);
\node[font=\tiny,gray!45!black,anchor=north] at (1.52,1.02) {$h_{d,3}$};
% toi phia truyen qua: vong qua mep be mat
\draw[gray!60,dashed,->] (1.05,2.25) .. controls (3.0,4.85) and (5.6,4.85) .. (6.60,3.72);
\node[font=\tiny,gray!45!black,anchor=south] at (4.40,4.62) {$h_{d,2}$};
\draw[gray!60,dashed,->] (1.05,1.45) .. controls (3.0,-1.05) and (5.6,-1.05) .. (6.80,0.30);
\node[font=\tiny,gray!45!black,anchor=north] at (4.40,-0.86) {$h_{d,4}$};

% --- nhan hai phia ---
\node[font=\tiny,cKenh,anchor=east] at (4.25,3.82) {phía phản xạ};
\node[font=\tiny,cKenh,anchor=west] at (4.55,3.82) {phía truyền qua};

\node[font=\tiny,align=left,anchor=north,text width=11cm] at (3.6,-1.55)
  {Đường liền: đường ghép tầng BS $\to$ STAR--RIS $\to$ UE.
   Đường nét đứt: đường trực tiếp BS $\to$ UE, tồn tại với mọi $k$.
   Kênh hiệu dụng của UE $k$ là $h^{\mathrm{eff}}_k=h_{d,k}+\sum_{n}h^{*}_{r,k,n}\,\phi^{X}_n\,g_n$
   với $X\in\{T,R\}$ theo phía của UE.};
\end{tikzpicture}
}
\caption{Mô hình hệ thống SISO STAR--RIS hỗ trợ RSMA}
\label{fig:ch2-mo-hinh}
\figuresource{Tác giả xây dựng theo mô hình trong mã nguồn và tham khảo \cite{Mu2021STARAided, Meng2024PPOSTARRSMA} (2026)}
\end{figure}

\section{Các giả định nghiên cứu}
\label{sec:gia-dinh}

Để tập trung vào bài toán phân bổ tài nguyên và so sánh thuật toán, luận văn sử dụng một
số giả định chuẩn hóa. Thứ nhất, thông tin trạng thái kênh (Channel State Information,
CSI) được biết tại thời điểm ra quyết định. Điều này cho phép tách bài toán ước lượng
kênh khỏi bài toán chọn hành động. Thứ hai, SIC được xem là lý tưởng; sau khi luồng chung
được giải mã thành công, phần tín hiệu tương ứng được loại bỏ hoàn toàn trước khi giải mã
luồng riêng. Thứ ba, STAR--RIS được mô hình hóa thụ động theo chế độ ES; công suất truyền
và phản xạ của mỗi phần tử có tổng bằng một. Thứ tư, pha truyền và pha phản xạ được phép
điều khiển độc lập trong mô hình hệ thống. Các ràng buộc ghép pha do phần cứng thực tế
chưa được xét và được nêu ở mục giới hạn nghiên cứu.

Kênh được sinh độc lập theo từng kịch bản. Mỗi hệ số là Gaussian phức không tương quan
theo mô hình trong tệp \texttt{physics.py}. Ba nhóm đường truyền có hệ số tỷ lệ khác nhau
để phản ánh mức độ mạnh tương đối: đường trực tiếp BS--UE dùng hệ số $0,35$, đường
BS--STAR--RIS dùng $1,0$ và đường STAR--RIS--UE dùng $0,75$. Các hệ số này là tham số mô
phỏng, không phải mô hình suy hao khoảng cách theo một bố trí địa lý cụ thể.

\section{Mô hình kênh}
\label{sec:mo-hinh-kenh}

\subsection{Sinh mẫu kênh phức}
\label{ssec:sinh-mau-kenh}

Một biến Gaussian phức chuẩn hóa trong mã nguồn được tạo từ hai biến Gaussian thực độc
lập. Ba khối kênh được mô tả như sau:
\begin{align}
h_{d,k} &\sim \mathcal{CN}(0,\ 0,35^2), \label{eq:kenh-truc-tiep}\\
g_n     &\sim \mathcal{CN}(0,\ 1^2),    \label{eq:kenh-bs-ris}\\
h_{r,k,n} &\sim \mathcal{CN}(0,\ 0,75^2). \label{eq:kenh-ris-ue}
\end{align}

Trong đó: $h_{d,k}$ là kênh trực tiếp từ BS đến UE $k$; $g_n$ là kênh từ BS đến phần tử
STAR--RIS thứ $n$; $h_{r,k,n}$ là kênh từ phần tử $n$ đến UE $k$; $\mathcal{CN}(0,
\sigma^2)$ là phân bố Gaussian phức tròn với trung bình bằng không và mức công suất tỷ lệ
với $\sigma^2$.

Ba công thức trên tạo ba chặng truyền độc lập. Đường trực tiếp được đặt yếu hơn đường đi
qua STAR--RIS để mô phỏng trường hợp bề mặt tái cấu hình có vai trò đáng kể. Sau đó, bộ
điều khiển không thay đổi các kênh gốc này mà thay đổi cách các thành phần qua STAR--RIS
cộng lại.

Để tiện trình bày, gom các hệ số BS--STAR--RIS thành véc-tơ $\mathbf{g} \in
\mathbb{C}^{N}$ và các hệ số STAR--RIS--UE $k$ thành véc-tơ $\mathbf{h}_{r,k} \in
\mathbb{C}^{N}$. Với $K$ UE, toàn bộ kênh STAR--RIS--UE có thể xem như ma trận kích thước
$K \times N$.

\subsection{Hệ số truyền và phản xạ của STAR--RIS}
\label{ssec:he-so-star-ris}

Với phần tử thứ $n$, hệ số phía truyền và phía phản xạ được viết:
\begin{equation}
\label{eq:phi-tr}
\phi^T_n = \sqrt{\beta^T_n}\, e^{j\theta^T_n},
\qquad
\phi^R_n = \sqrt{1 - \beta^T_n}\, e^{j\theta^R_n}.
\end{equation}

Trong đó: $\beta^T_n \in [0,1]$ là tỷ lệ năng lượng dành cho phía truyền; $1 - \beta^T_n$
là tỷ lệ dành cho phía phản xạ; $\theta^T_n, \theta^R_n \in [-\pi, \pi)$ là hai pha điều
khiển; $j$ là đơn vị ảo. Công thức cho thấy STAR--RIS có hai loại nút điều khiển trên mỗi
phần tử: biến $\beta$ giống một van chia năng lượng giữa hai phía, còn $\theta$ giống một
núm xoay pha dùng để làm các đường truyền cộng cùng hướng tại UE.

Đặt hai ma trận đường chéo:
\begin{equation}
\label{eq:ma-tran-duong-cheo}
\boldsymbol{\Phi}^T = \mathrm{diag}(\phi^T_1, \dots, \phi^T_N),
\qquad
\boldsymbol{\Phi}^R = \mathrm{diag}(\phi^R_1, \dots, \phi^R_N).
\end{equation}

Trong đó: $\mathrm{diag}(\cdot)$ tạo ma trận đường chéo và mỗi phần tử đường chéo chứa hệ
số truyền hoặc phản xạ của một phần tử STAR--RIS. Đây chỉ là cách gom $N$ phần tử riêng
lẻ vào một đối tượng toán học; về vật lý, mỗi phần tử vẫn tác động độc lập lên tín hiệu
đi qua vị trí của nó.

\subsection{Kênh hiệu dụng}
\label{ssec:kenh-hieu-dung}

Với UE $k$, kênh hiệu dụng gồm đường trực tiếp cộng với đường ghép tầng qua STAR--RIS:
\begin{equation}
\label{eq:kenh-hieu-dung}
h^{\mathrm{eff}}_k = h_{d,k} + \mathbf{h}^{H}_{r,k}\,\boldsymbol{\Phi}^{s_k}\,\mathbf{g},
\qquad s_k \in \{T, R\}.
\end{equation}

Trong đó: $h^{\mathrm{eff}}_k$ là kênh tổng mà UE $k$ thực sự nhìn thấy; $h_{d,k}$ là
đường trực tiếp; $\mathbf{h}^{H}_{r,k}$ là liên hợp chuyển vị của kênh STAR--RIS--UE;
$\boldsymbol{\Phi}^{s_k}$ chọn ma trận truyền hoặc phản xạ theo phía của UE; $\mathbf{g}$
là kênh BS--STAR--RIS.

Công thức~\eqref{eq:kenh-hieu-dung} là cầu nối giữa STAR--RIS và RSMA. Bộ điều khiển
không thay đổi $h_{d,k}$, $\mathbf{g}$ hay $\mathbf{h}_{r,k}$; nó thay đổi $\beta$ và
$\theta$ trong $\boldsymbol{\Phi}$, từ đó làm đường ghép tầng cộng thuận hoặc nghịch với
đường trực tiếp. Giá trị cuối cùng $|h^{\mathrm{eff}}_k|$ quyết định mức mạnh tín hiệu
trong các công thức tỷ số tín hiệu trên nhiễu cộng tạp âm phía sau. Mã nguồn tính trực
tiếp tổng từng phần tử thay vì nhân ma trận để giảm chi phí và giữ biểu thức rõ ràng; hai
cách viết tương đương về mặt toán học.

\section{Mô hình tín hiệu RSMA}
\label{sec:mo-hinh-tin-hieu}

\subsection{Tín hiệu phát tại trạm gốc}
\label{ssec:tin-hieu-phat}

RSMA một lớp tạo một ký hiệu luồng chung $s_c$ và $K$ ký hiệu luồng riêng $s_k$. Tín hiệu
phát của BS là:
\begin{equation}
\label{eq:tin-hieu-phat}
x = \sqrt{p_c}\, s_c + \sum_{k=1}^{K} \sqrt{p_k}\, s_k .
\end{equation}

Trong đó: $p_c$ là công suất luồng chung; $p_k$ là công suất luồng riêng của UE $k$; $s_c$
là ký hiệu chung; $s_k$ là ký hiệu riêng; các ký hiệu được chuẩn hóa công suất đơn vị.

Công thức~\eqref{eq:tin-hieu-phat} cho thấy BS chỉ có một anten và phát một tín hiệu tổng
hợp duy nhất. Không có véc-tơ tạo búp sóng riêng cho từng UE. Vì vậy, ở một UE bất kỳ, tất
cả luồng riêng đều đi qua cùng một hệ số kênh vô hướng $h^{\mathrm{eff}}_k$. Đây chính là
nguyên nhân làm nhiễu luồng riêng trở nên đáng chú ý trong hệ SISO.

Ràng buộc tổng công suất được viết:
\begin{equation}
\label{eq:rang-buoc-cong-suat}
p_c + \sum_{k=1}^{K} p_k = P_{\max},
\qquad p_c \ge 0,\ p_k \ge 0.
\end{equation}

Trong đó: $P_{\max}$ là ngân sách công suất phát chuẩn hóa của BS; các biến công suất
không âm và tổng bằng ngân sách. Công thức thể hiện quy tắc chia một chiếc bánh công
suất: tăng công suất cho luồng này đồng nghĩa giảm phần còn lại cho các luồng khác. Bộ
giải mã hành động ở Chương~\ref{chap:phuong-phap} được thiết kế để ràng buộc này luôn
đúng mà không cần phạt sau khi hành động đã được sinh.

Hình~\ref{fig:ch2-kien-truc} trình bày kiến trúc thu--phát tương ứng. Ở phía phát, thông
điệp của từng UE đi qua bộ tách thông điệp, phần chung được gộp lại thành một luồng, rồi
toàn bộ đi qua bộ mã hóa và khối phân bổ công suất trước khi phát qua một anten duy nhất.
Ở phía thu, mỗi UE giải mã luồng chung, dùng một tầng SIC để loại luồng chung, giải mã
luồng riêng của mình rồi ghép hai phần lại.

\begin{figure}[htbp]
\centering
\resizebox{\textwidth}{!}{\begin{tikzpicture}[font=\scriptsize,>=Stealth,line width=0.55pt]

% ============ PHIA PHAT (BS SISO) ============
\node[vbox,text width=2.6cm] (sch) at (0,0.6) {Lập lịch $K$ UE};
\node[vbox,text width=3.2cm] (spl) at (1.6,0.6) {Bộ tách thông điệp};
\node[vbox,text width=1.5cm] (cmb) at (3.1,1.50) {Bộ ghép\\ phần chung};
\node[vbox,text width=3.2cm] (enc) at (4.6,0.6) {Bộ mã hóa};
\node[vbox,text width=3.2cm] (pwr) at (6.1,0.6) {Phân bổ công suất $p_c,\,p_k$};

% thong diep vao
\draw[->] ($(sch.south)+(0,0.85)$)  -- node[above,inner sep=1.5pt]{\lbm{W_1}}   ($(spl.north)+(0,0.85)$);
\draw[->] ($(sch.south)+(0,0)$)     -- node[above,inner sep=1.5pt]{\lbm{W_k}}   ($(spl.north)+(0,0)$);
\draw[->] ($(sch.south)+(0,-0.85)$) -- node[above,inner sep=1.5pt]{\lbm{W_K}}   ($(spl.north)+(0,-0.85)$);
\node[font=\tiny] at (0.8,0.30) {$\vdots$};
\node[font=\tiny] at (0.8,-0.55) {$\vdots$};

% phan chung -> bo ghep
\draw[->] (1.88,2.05) -- node[above,inner sep=1pt,font=\tiny]{\lbm{W_{c,1}}} (2.65,2.05);
\draw[->] (1.88,1.50) -- node[above,inner sep=1pt,font=\tiny]{\lbm{W_{c,k}}} (2.65,1.50);
\draw[->] (1.88,0.95) -- node[above,inner sep=1pt,font=\tiny]{\lbm{W_{c,K}}} (2.65,0.95);

% phan rieng -> bo ma hoa
\draw[->] (1.88,0.10)  -- node[above,inner sep=1pt,font=\tiny]{\lbm{W_{p,1}}} (4.32,0.10);
\draw[->] (1.88,-0.40) -- node[above,inner sep=1pt,font=\tiny]{\lbm{W_{p,k}}} (4.32,-0.40);
\draw[->] (1.88,-0.90) -- node[above,inner sep=1pt,font=\tiny]{\lbm{W_{p,K}}} (4.32,-0.90);

% bo ghep -> bo ma hoa
\draw[->] (3.55,1.50) -- node[above,inner sep=1pt,font=\tiny]{\textcolor{magenta}{$W_c$}} (4.32,1.50);

% bo ma hoa -> phan bo cong suat
\draw[->] (4.88,1.50)  -- node[above,inner sep=1pt,font=\tiny]{\textcolor{magenta}{$s_c$}} (5.82,1.50);
\draw[->] (4.88,0.10)  -- node[above,inner sep=1pt,font=\tiny]{$s_1$} (5.82,0.10);
\draw[->] (4.88,-0.40) -- node[above,inner sep=1pt,font=\tiny]{$s_k$} (5.82,-0.40);
\draw[->] (4.88,-0.90) -- node[above,inner sep=1pt,font=\tiny]{$s_K$} (5.82,-0.90);

% mot anten phat duy nhat (SISO)
\path (8.0,0.6) pic {ant};
\draw (pwr.south) -- (8.0,0.6);
\node[anchor=south,inner sep=1pt] at (7.55,0.72) {\sgm{x}};
\node[font=\tiny,anchor=north] at (6.55,-1.85)
      {$x=\sqrt{p_c}\,s_c+\sum_{k=1}^{K}\sqrt{p_k}\,s_k$};

% ============ UE 1 ============
\path (10.4,3.00) pic[rotate=180] {ant};
\node[anchor=south west,inner sep=1pt] at (10.28,3.12) {\sgm{y_1}};
\node[draw=blue,line width=0.9pt,fill=blue!8,rounded corners=2pt,inner sep=4pt] (ue1) at (11.75,3.00) {UE$_1$};
\draw[->] (10.4,3.00) -- (ue1.west);
\draw[->] (ue1.east) -- (13.6,3.00) node[midway,above,inner sep=1.5pt]{\lbm{\widehat{W}_1}};
\node[font=\tiny] at (10.4,2.35) {$\vdots$};

% ============ UE k (chi tiet) ============
\path (10.4,1.85) pic[rotate=180] {ant};
\node[anchor=south east,inner sep=1pt] at (10.42,1.99) {\sgm{y_k}};
\node[box,text width=1.35cm] (deca) at (12.15,1.85) {Giải mã\\ luồng chung};
\node[box] (spl2) at (14.35,1.85) {Tách};
\node[sbox] (encs) at (14.35,0.85) {Mã hóa};
\node[sbox] (chs)  at (12.60,0.85) {Kênh};
\node[draw,circle,inner sep=1.2pt,line width=0.4pt] (sum) at (11.40,0.85) {\tiny$+$};
\node[font=\tiny] at (11.62,1.10) {SIC};
\node[box,text width=1.35cm] (decb) at (12.45,-0.40) {Giải mã\\ luồng riêng};
\node[vbox,text width=2.4cm,minimum height=0.5cm] (cmb2) at (15.95,0.60) {Ghép};

\draw[->] (10.4,1.85) -- (deca.west);
\draw[->] (deca.east) -- node[above,inner sep=1.5pt]{\textcolor{magenta}{$\widehat{W}_c$}} (spl2.west);
\draw[->] (spl2.east) -- (15.55,1.85) node[midway,above,inner sep=1pt,font=\tiny]{$\widehat{W}_{c,k}$}
          -- (15.55,1.05) -- (15.70,1.05);
\draw[->] (spl2.south) -- (encs.north);
\draw[->] (encs.west) -- (chs.east);
\draw[->] (chs.west) -- (sum.east);
\draw[->] (10.95,1.85) -- (10.95,0.85) -- (sum.west);
\node[font=\tiny] at (11.20,0.99) {$-$};
\draw[->] (sum.south) -- (10.95,-0.40) -- (decb.west);
\draw[->] (decb.east) -- (15.55,-0.40) node[midway,above,inner sep=1pt,font=\tiny]{$\widehat{W}_{p,k}$}
          -- (15.55,0.15) -- (15.70,0.15);
\draw[->] (cmb2.south) -- (17.45,0.60);
\node[anchor=south,inner sep=2pt] at (16.95,0.68) {\lbm{\widehat{W}_k}};

% ============ UE K ============
\path (10.4,-1.80) pic[rotate=180] {ant};
\node[anchor=south west,inner sep=1pt] at (10.28,-1.68) {\sgm{y_K}};
\node[draw=violet,line width=0.9pt,fill=violet!8,rounded corners=2pt,inner sep=4pt] (ueK) at (11.75,-1.80) {UE$_K$};
\draw[->] (10.4,-1.80) -- (ueK.west);
\draw[->] (ueK.east) -- (13.6,-1.80) node[midway,above,inner sep=1.5pt]{\lbm{\widehat{W}_K}};

% ============ NEN: kenh, cac khung nhom ============
\begin{scope}[on background layer]
  \node[cloud,cloud puffs=15,draw=green!45!black,fill=green!16,
        minimum width=3.2cm,minimum height=6.4cm] at (9.30,0.60) {};
  \node[align=center,font=\scriptsize] at (9.25,0.75)
       {Kênh hiệu dụng\\[1pt] $h^{\mathrm{eff}}_k$\\[2pt] \tiny (đường trực tiếp\\[-1pt] \tiny + STAR--RIS)};
  \grp{1.05}{-1.55}{6.65}{2.45}{magenta}{Trạm gốc BS (SISO)}
  \grp{10.78}{-1.05}{16.55}{2.42}{orange}{UE$_k$}
  \node[draw=black!55,line width=0.4pt,inner sep=0pt,
        fit={(11.15,0.45)(15.15,1.28)}] {};
\end{scope}
\node[inner sep=0pt,fit={(-0.5,-2.75)(17.95,3.75)}] {};  % giu khung bao, khong ve vien
\end{tikzpicture}
}
\caption{Kiến trúc thu--phát RSMA một lớp cho $K$ UE trong hệ SISO}
\label{fig:ch2-kien-truc}
\figuresource{Tác giả xây dựng dựa trên \cite{Mao2018RSMABridging, Mao2022RSMAFundamentals} (2026)}
\end{figure}

\subsection{Tín hiệu thu}
\label{ssec:tin-hieu-thu}

Tín hiệu thu tại UE $k$ là:
\begin{equation}
\label{eq:tin-hieu-thu}
y_k = h^{\mathrm{eff}}_k\, x + n_k .
\end{equation}

Trong đó: $y_k$ là tín hiệu UE $k$ nhận được; $h^{\mathrm{eff}}_k$ là kênh hiệu dụng; $x$
là tín hiệu tổng hợp của BS; $n_k \sim \mathcal{CN}(0, \sigma^2)$ là tạp âm thu. Công
thức có ý nghĩa rất trực tiếp: mọi thành phần trong $x$ đều bị nhân bởi cùng một kênh vô
hướng của UE $k$. Do đó, nếu nhiều luồng riêng cùng bật, UE không có cơ chế không gian để
làm yếu riêng từng luồng nhiễu; khác biệt chỉ đến từ công suất của các luồng.

\subsection{Giải mã luồng chung}
\label{ssec:giai-ma-chung}

Ở bước đầu, UE $k$ giải mã luồng chung và xem toàn bộ luồng riêng như nhiễu. Tỷ số tín
hiệu trên nhiễu cộng tạp âm (Signal-to-Interference-plus-Noise Ratio, SINR) của luồng
chung là:
\begin{equation}
\label{eq:sinr-chung}
\gamma_{c,k} = \frac{|h^{\mathrm{eff}}_k|^2 p_c}
                    {|h^{\mathrm{eff}}_k|^2 \sum_{j=1}^{K} p_j + \sigma^2}.
\end{equation}

Trong đó: $\gamma_{c,k}$ là SINR khi UE $k$ giải mã luồng chung; tử số là công suất luồng
chung sau kênh; mẫu số gồm toàn bộ công suất luồng riêng sau cùng kênh và công suất tạp
âm $\sigma^2$. Công thức cho thấy luồng chung muốn dễ giải mã cần có tỷ lệ công suất đủ
lớn so với tổng công suất luồng riêng. Đồng thời, UE có kênh yếu sẽ chịu ảnh hưởng tạp âm
lớn hơn và có thể trở thành nút cổ chai của luồng chung.

Tốc độ luồng chung mà UE $k$ có thể giải mã là:
\begin{equation}
\label{eq:toc-do-chung-k}
R_{c,k} = \log_2(1 + \gamma_{c,k}).
\end{equation}

Trong đó: $R_{c,k}$ có đơn vị bit/s/Hz và biểu diễn tốc độ tối đa của luồng chung tại UE
$k$ theo công thức Shannon với SINR tương ứng. SINR tăng theo tỷ lệ nhưng tốc độ tăng
theo hàm logarit, do đó lợi ích của việc tiếp tục tăng công suất có xu hướng giảm dần.

Vì tất cả UE phải giải mã cùng một luồng chung, tốc độ chung của hệ thống phải không vượt
khả năng của UE yếu nhất:
\begin{equation}
\label{eq:toc-do-chung}
R_c = \min_{k \in \{1,\dots,K\}} R_{c,k}.
\end{equation}

Trong đó: $R_c$ là tốc độ chung dùng cho toàn hệ thống và phép lấy nhỏ nhất chọn UE có
khả năng giải mã thấp nhất.

Công thức~\eqref{eq:toc-do-chung} giải thích vì sao căn pha chỉ cho UE mạnh nhất chưa đủ.
Nếu STAR--RIS làm kênh của một UE khác quá yếu, UE đó kéo $R_c$ xuống và làm giảm phần
tốc độ chung của tất cả người dùng. Bộ điều khiển vì vậy phải cân bằng giữa tăng độ lợi
cho UE được ưu tiên và giữ khả năng giải mã luồng chung của các UE còn lại.

Sau khi $R_c$ được xác định, tốc độ chung được chia cho từng UE thông qua véc-tơ $\eta$:
\begin{equation}
\label{eq:phan-bo-chung}
C_k = \eta_k R_c,
\qquad \eta_k \ge 0,
\qquad \sum_{k=1}^{K} \eta_k = 1.
\end{equation}

Trong đó: $C_k$ là phần tốc độ chung cấp cho UE $k$; $\eta_k$ là tỷ lệ phân bổ; tổng các
tỷ lệ bằng một nên toàn bộ $R_c$ được phân phối hết. Công thức biến luồng chung thành một
quỹ tốc độ có thể phân bổ linh hoạt: UE có luồng riêng yếu có thể nhận $C_k$ lớn hơn để
đạt QoS, còn UE đã có tốc độ riêng cao có thể nhận ít hơn. Vì vậy, $\eta$ là biến tài
nguyên thực sự chứ không chỉ là ký hiệu phụ.

\subsection{Khử nhiễu liên tiếp và giải mã luồng riêng}
\label{ssec:sic-giai-ma-rieng}

Sau khi giải mã $s_c$, UE loại luồng chung khỏi tín hiệu bằng SIC. Khi đó, UE $k$ giải mã
$s_k$ và xem các luồng riêng còn lại như nhiễu. SINR luồng riêng là:
\begin{equation}
\label{eq:sinr-rieng}
\gamma_{p,k} = \frac{|h^{\mathrm{eff}}_k|^2 p_k}
                    {|h^{\mathrm{eff}}_k|^2 \sum_{j \neq k} p_j + \sigma^2}.
\end{equation}

Trong đó: $\gamma_{p,k}$ là SINR của luồng riêng UE $k$; tử số là công suất của chính
luồng $k$; tổng trong mẫu số chứa toàn bộ công suất luồng riêng của các UE khác; luồng
chung không còn trong mẫu vì đã được SIC loại bỏ.

Công thức~\eqref{eq:sinr-rieng} là phần lõi của phân tích SISO. Cùng một độ lợi kênh
$|h^{\mathrm{eff}}_k|^2$ nhân cả tín hiệu mong muốn và các luồng riêng gây nhiễu. Khi tạp
âm nhỏ, độ lợi kênh gần như triệt tiêu khỏi tỷ số, nên tăng toàn bộ mức kênh không giải
quyết được nhiễu giữa các luồng riêng.

Tốc độ luồng riêng của UE $k$ là:
\begin{equation}
\label{eq:toc-do-rieng}
R_{p,k} = \log_2(1 + \gamma_{p,k}).
\end{equation}

Trong đó: $R_{p,k}$ là tốc độ do luồng riêng mang đến cho UE $k$. Nếu SINR luồng riêng bị
đóng trần bởi nhiễu giữa các luồng riêng thì tốc độ riêng cũng bị giới hạn, dù độ lợi
kênh tổng thể tăng mạnh. Đây là lý do phân bổ công suất luồng riêng gần một luồng chiếm
ưu thế có thể hiệu quả trong hệ SISO.

Tốc độ cuối cùng của UE và tổng tốc độ hệ thống là:
\begin{align}
R_k &= C_k + R_{p,k}, \label{eq:toc-do-ue}\\
R_{\mathrm{sum}} &= \sum_{k=1}^{K} R_k. \label{eq:tong-toc-do}
\end{align}

Trong đó: $R_k$ là tổng tốc độ của UE $k$ gồm phần chung và phần riêng;
$R_{\mathrm{sum}}$ là tổng tốc độ của toàn bộ $K$ UE. Hai công thức cho thấy RSMA có hai
con đường để cấp tốc độ cho một UE: tăng $R_{p,k}$ bằng công suất riêng và kênh thuận
lợi, hoặc tăng $C_k$ bằng cách phân bổ phần luồng chung. Chính sự linh hoạt này giúp đồng
thời theo đuổi tổng tốc độ và QoS.

\section{Chỉ số chất lượng dịch vụ và hàm đánh giá}
\label{sec:qos}

Luận văn dùng tốc độ tối thiểu $R_{\min}$ làm yêu cầu QoS. Với từng UE, phần thiếu hụt
tốc độ được định nghĩa:
\begin{equation}
\label{eq:thieu-hut}
d_k = \max(R_{\min} - R_k,\ 0).
\end{equation}

Trong đó: $d_k$ là mức thiếu hụt tốc độ của UE $k$; nếu UE đạt hoặc vượt $R_{\min}$ thì
$d_k = 0$. Công thức biến một ràng buộc đạt hoặc không đạt thành đại lượng có độ lớn. UE
thiếu $0,01$ bit/s/Hz bị xem nhẹ hơn UE thiếu $0,4$ bit/s/Hz, nhờ đó thuật toán học nhận
được tín hiệu điều chỉnh mượt hơn.

Hai chỉ số vi phạm dùng trong mã nguồn là:
\begin{equation}
\label{eq:vi-pham}
V = \sum_{k=1}^{K} d_k,
\qquad
V_2 = \sum_{k=1}^{K} d_k^2 .
\end{equation}

Trong đó: $V$ là tổng thiếu hụt tuyến tính và $V_2$ là tổng bình phương thiếu hụt. Hai
đại lượng kết hợp hai góc nhìn: thành phần tuyến tính phản ứng với mọi vi phạm, còn thành
phần bình phương phạt mạnh hơn khi một UE thiếu QoS nhiều. Điều này giúp phần thưởng
tránh ưu tiên một nghiệm có tổng tốc độ cao nhưng bỏ rơi một UE quá mức.

Ngoài ra, hai chỉ số dễ diễn giải được báo cáo: tỷ lệ UE đạt QoS trong một kịch bản và
biến nhị phân cho biết tất cả UE có đạt QoS hay không. Các chỉ số này không trực tiếp tối
ưu nhưng dùng để kiểm tra chất lượng nghiệm bên cạnh tổng tốc độ.

\section{Các biến quyết định và ràng buộc}
\label{sec:bien-rang-buoc}

Toàn bộ hành động vật lý có thể nhóm thành:
\begin{equation}
\label{eq:tap-hanh-dong}
\mathcal{A} = \{\, \mathbf{p},\ \boldsymbol{\eta},\ \boldsymbol{\beta},\
\boldsymbol{\theta}^T,\ \boldsymbol{\theta}^R \,\},
\qquad \mathbf{p} = [p_c, p_1, \dots, p_K].
\end{equation}

Trong đó: $\mathbf{p}$ là véc-tơ công suất; $\boldsymbol{\eta}$ là phân bổ tốc độ chung;
$\boldsymbol{\beta}$ là hệ số chia năng lượng phía truyền; $\boldsymbol{\theta}^T$ và
$\boldsymbol{\theta}^R$ là các pha truyền và phản xạ. Một quyết định tài nguyên vì vậy
không phải một số duy nhất mà là một véc-tơ lớn gồm các biến tác động ở nhiều tầng: RSMA
quyết định phát gì và chia công suất ra sao, còn STAR--RIS quyết định kênh hiệu dụng được
tạo ra như thế nào.

Các ràng buộc chính được viết gọn:
\begin{align}
& p_i \ge 0, \qquad \sum_{i=0}^{K} p_i = P_{\max}, \label{eq:rb-cong-suat}\\
& \eta_k \ge 0, \qquad \sum_{k=1}^{K} \eta_k = 1,  \label{eq:rb-eta}\\
& 0 \le \beta^T_n \le 1,                            \label{eq:rb-beta}\\
& -\pi \le \theta^T_n,\ \theta^R_n < \pi,           \label{eq:rb-pha}\\
& R_k \ge R_{\min}, \quad \forall k.                \label{eq:rb-qos}
\end{align}

Trong đó: $i = 0$ tương ứng luồng chung; $k$ đánh số UE; $n$ đánh số phần tử STAR--RIS;
các dòng lần lượt là ràng buộc công suất, ràng buộc phân bổ tốc độ chung, ràng buộc chia
năng lượng, miền pha và QoS.

Năm dòng trên mô tả biên vật lý của bài toán. Bốn dòng đầu có thể được bảo đảm trực tiếp
bằng cách thiết kế bộ giải mã phù hợp. Dòng QoS phức tạp hơn vì $R_k$ phụ thuộc đồng thời
vào tất cả công suất và biến STAR--RIS, nên luận văn dùng cả tiêu chí xác thực và hạng
phạt trong phần thưởng để điều khiển.

\section{Phát biểu bài toán tối ưu}
\label{sec:phat-bieu-bai-toan}

Bài toán tham chiếu của luận văn là:
\begin{align}
\max_{\mathbf{p},\,\boldsymbol{\eta},\,\boldsymbol{\beta},\,
      \boldsymbol{\theta}^T,\,\boldsymbol{\theta}^R} \quad & R_{\mathrm{sum}}
      \label{eq:bai-toan}\\
\text{với điều kiện} \quad & \eqref{eq:rb-cong-suat}\text{--}\eqref{eq:rb-qos}.
      \label{eq:bai-toan-dk}
\end{align}

Trong đó: biến tối ưu gồm toàn bộ năm nhóm tài nguyên; hàm mục tiêu là tổng tốc độ; các
điều kiện là các ràng buộc vật lý và QoS đã nêu. Đây là mục tiêu cuối cùng mà mọi phương
pháp đều hướng tới. AO cố cải thiện trực tiếp nghiệm theo từng CSI, còn DRL không giải
lại bài toán từ đầu ở thời điểm sử dụng mà học một ánh xạ từ CSI sang hành động trong
giai đoạn huấn luyện.

Bài toán có tính phi lồi vì nhiều nguyên nhân. Pha xuất hiện bên trong tổng số phức và
sau đó được lấy trị tuyệt đối bình phương; các biến công suất xuất hiện ở cả tử và mẫu
của SINR; tốc độ chung có phép lấy nhỏ nhất trên UE; QoS ghép tất cả biến qua tốc độ.
Ngoài ra, khi $N$ tăng, số biến pha và năng lượng tăng tuyến tính làm miền tìm kiếm lớn
nhanh.

Số chiều hành động vật lý được dùng trong mã nguồn là:
\begin{equation}
\label{eq:so-chieu-hanh-dong}
d_a = (K + 1) + K + 3N = 2K + 1 + 3N.
\end{equation}

Trong đó: $K+1$ biến đầu là công suất luồng chung và riêng; $K$ biến tiếp theo là phân bổ
tốc độ chung; $3N$ biến cuối gồm $N$ hệ số năng lượng, $N$ pha truyền và $N$ pha phản xạ.
Độ khó vì vậy tăng theo kích thước STAR--RIS: với $K = 4$ và $N = 32$, mạng phải sinh
$105$ giá trị hành động; với $N = 128$, số chiều tăng lên $393$. Nếu mỗi giá trị được học
hoàn toàn tự do, việc khám phá các vùng hành động tốt trở nên khó hơn.

\section{Phân tích cấu trúc SISO--RSMA}
\label{sec:phan-tich-cau-truc}

\subsection{Nhiễu luồng riêng trong miền tỷ số tín hiệu trên tạp âm cao}
\label{ssec:nhieu-luong-rieng}

Đặt $g_k = |h^{\mathrm{eff}}_k|^2$ và $P_{\mathrm{priv}} = \sum_j p_j$. Công thức SINR
luồng riêng có thể viết lại:
\begin{equation}
\label{eq:sinr-rieng-vietlai}
\gamma_{p,k} = \frac{g_k p_k}{g_k (P_{\mathrm{priv}} - p_k) + \sigma^2}.
\end{equation}

Trong đó: $g_k$ là độ lợi kênh hiệu dụng; $P_{\mathrm{priv}}$ là tổng công suất của tất
cả luồng riêng; $P_{\mathrm{priv}} - p_k$ là tổng công suất các luồng riêng gây nhiễu cho
UE $k$. Công thức làm lộ rõ đặc điểm SISO: cùng một $g_k$ nhân cả tín hiệu mong muốn và
nhiễu riêng, vì vậy khi nhiễu riêng mạnh hơn tạp âm, lợi ích của việc tăng độ lợi kênh
đối với tỷ số này giảm đáng kể.

Khi $g_k (P_{\mathrm{priv}} - p_k) \gg \sigma^2$, có xấp xỉ:
\begin{equation}
\label{eq:xap-xi-nhieu}
\gamma_{p,k} \approx \frac{p_k}{P_{\mathrm{priv}} - p_k}.
\end{equation}

Trong đó: điều kiện xấp xỉ nghĩa là hệ thống đang ở vùng nhiễu luồng riêng chi phối tạp
âm; khi đó $g_k$ gần như triệt tiêu khỏi tử và mẫu. Công thức giải thích vì sao chỉ làm
kênh mạnh hơn chưa chắc làm tốc độ riêng tăng nếu nhiều luồng riêng cùng hoạt động: yếu
tố quyết định chuyển thành tỷ lệ công suất giữa luồng mong muốn và các luồng riêng khác.

Nếu có $M \ge 2$ luồng riêng chia đều công suất, mỗi luồng có $p_k = P_{\mathrm{priv}}/M$.
Khi đó:
\begin{equation}
\label{eq:tran-toc-do}
\gamma_{p,k} \approx \frac{1}{M-1},
\qquad
R_{p,k} \approx \log_2\!\left(\frac{M}{M-1}\right).
\end{equation}

Trong đó: $M$ là số luồng riêng đang hoạt động với mức công suất bằng nhau. Công thức cho
một trực giác rất rõ: với hai luồng riêng chia đều, mỗi luồng chỉ đạt khoảng $1$ bit/s/Hz
trong miền nhiễu chi phối; với bốn luồng chia đều, mỗi luồng khoảng $\log_2(4/3) \approx
0,415$ bit/s/Hz. Các giá trị này gần như không phụ thuộc vào kênh khi xấp xỉ đúng.

Ngược lại, nếu chỉ một luồng riêng $k^\star$ mang gần toàn bộ $P_{\mathrm{priv}}$, mẫu số
của luồng đó gần như chỉ còn tạp âm:
\begin{equation}
\label{eq:tap-trung-cong-suat}
\gamma_{p,k^\star} \approx \frac{g_{k^\star} P_{\mathrm{priv}}}{\sigma^2}.
\end{equation}

Trong đó: $k^\star$ là UE được ưu tiên luồng riêng và các luồng riêng còn lại có công
suất rất nhỏ. Khi giảm rò rỉ công suất sang các luồng riêng khác, UE được chọn thoát khỏi
giới hạn do nhiễu liên luồng và tốc độ có thể tăng theo tỷ số tín hiệu trên tạp âm. Đây
là cơ sở vật lý cho việc dùng softmax có độ sắc cao trong bộ giải mã công suất riêng.

Hình~\ref{fig:ch2-tran-toc-do} minh họa hai kết quả trên. Khung bên trái vẽ $R_{p,k}$
theo tỷ số tín hiệu trên tạp âm: đường ứng với $M = 1$ tăng gần tuyến tính theo thang
decibel, trong khi hai đường ứng với $M = 2$ và $M = 4$ bão hòa ngay từ khoảng mười
decibel và không vượt được trần của mình dù kênh có tốt đến đâu. Khung bên phải vẽ giá
trị tiệm cận $\log_2\big(M/(M-1)\big)$ theo số luồng riêng hoạt động, cho thấy trần này
giảm rất nhanh khi $M$ tăng.

\begin{figure}[htbp]
\centering
\resizebox{\textwidth}{!}{\begin{tikzpicture}[font=\footnotesize]
\begin{axis}[
  name=left, width=8.2cm, height=6.0cm,
  xlabel={Tỷ số tín hiệu trên tạp âm $g_kP_{\mathrm{priv}}/\sigma^2$ (dB)},
  ylabel={$R_{p,k}$ (bit/s/Hz)},
  xmin=0,xmax=30, ymin=0, ymax=10.5, ytick={0,2,4,6,8,10},
  grid=both, grid style={gray!25}, tick label style={font=\scriptsize},
  label style={font=\scriptsize},
  legend style={font=\scriptsize,at={(0.03,0.97)},anchor=north west,draw=gray!50},
  legend cell align=left,
]
% M = 1 : toan bo cong suat rieng don vao mot luong
\addplot[very thick,blue!70!black,domain=0:30,samples=120]
   {log2(1 + 10^(x/10))};
\addlegendentry{$M=1$ (tập trung)}
% M = 2
\addplot[very thick,orange!85!black,dashed,domain=0:30,samples=120]
   {log2(1 + 10^(x/10)/(1*10^(x/10) + 2))};
\addlegendentry{$M=2$ (chia đều)}
% M = 4
\addplot[very thick,violet,dotted,domain=0:30,samples=120]
   {log2(1 + 10^(x/10)/(3*10^(x/10) + 4))};
\addlegendentry{$M=4$ (chia đều)}
% duong tiem can
\addplot[gray,thin,domain=0:30,samples=2]{1.0};
\addplot[gray,thin,domain=0:30,samples=2]{0.415};
\node[font=\tiny,anchor=west,gray!45!black,align=left] at (axis cs:13.5,1.9)
     {đường xám: trần tiệm cận $\log_2\!\big(M/(M-1)\big)$};
\end{axis}

\begin{axis}[
  at={(left.east)}, anchor=west, xshift=1.1cm,
  width=6.6cm, height=6.0cm,
  xlabel={Số luồng riêng hoạt động $M$},
  ylabel={$R_{p,k}$ tiệm cận (bit/s/Hz)},
  ybar, bar width=9pt, xmin=1.4,xmax=8.6, ymin=0,ymax=1.2,
  xtick={2,3,4,5,6,7,8},
  grid=both, grid style={gray!25}, tick label style={font=\scriptsize},
  label style={font=\scriptsize},
  nodes near coords, every node near coord/.append style={font=\tiny,/pgf/number format/.cd,fixed,precision=3,use comma},
]
\addplot[fill=blue!35,draw=blue!60!black] coordinates
  {(2,1.000)(3,0.585)(4,0.415)(5,0.322)(6,0.263)(7,0.222)(8,0.193)};
\end{axis}
\end{tikzpicture}
}
\caption{Trần tốc độ luồng riêng do nhiễu giữa các luồng riêng trong hệ SISO}
\label{fig:ch2-tran-toc-do}
\figuresource{Tác giả xây dựng từ công thức \eqref{eq:xap-xi-nhieu}--\eqref{eq:tap-trung-cong-suat} (2026)}
\end{figure}

Phân tích trên không phải chứng minh rằng nghiệm tối ưu toàn cục luôn dồn toàn bộ công suất
riêng vào đúng một luồng, tức dạng one-hot theo cách gọi trong học máy.
Khi QoS, luồng chung, tỷ số tín hiệu trên tạp âm hữu hạn và sự thay đổi của kênh được xét
đầy đủ, một số kịch bản có thể hưởng lợi từ nhiều luồng riêng. Kết luận an toàn là bài
toán có xu hướng ưu tiên phân bổ riêng tập trung trong nhiều trạng thái kênh SISO, vì vậy
một tham số hóa có khả năng biểu diễn phân bổ gần one-hot là hợp lý.

\subsection{Vai trò của luồng chung đối với chất lượng dịch vụ}
\label{ssec:vai-tro-luong-chung}

Giả sử một UE không được cấp công suất riêng đáng kể, tốc độ của nó chủ yếu đến từ
$C_k = \eta_k R_c$. Như vậy, luồng chung có thể đóng vai trò mức nền dịch vụ. Nếu $R_c$
đủ lớn và $\eta$ được phân bổ hợp lý, các UE không được chọn làm luồng riêng chính vẫn có
thể đạt $R_{\min}$.

Trong trường hợp đơn giản phân bổ đều tốc độ chung cho $K = 4$ UE, yêu cầu $C_k \ge 0,5$
bit/s/Hz dẫn đến $R_c \ge 2$ bit/s/Hz. Tuy nhiên, đây chỉ là ví dụ minh họa. Trong mô
hình thực, $\eta$ được tối ưu nên UE đã có tốc độ riêng cao không cần nhận phần tốc độ
chung lớn. Vì vậy, ngưỡng công suất luồng chung không thể suy ra chỉ từ phân bổ đều mà
phải xem là một lựa chọn thiết kế được kiểm chứng bằng thực nghiệm.

\subsection{Quan hệ công suất chung và riêng trong miền tỷ số tín hiệu trên tạp âm cao}
\label{ssec:quan-he-chung-rieng}

Khi một luồng riêng chiếm ưu thế và UE nút cổ chai của luồng chung vẫn ở vùng tỷ số tín
hiệu trên tạp âm cao, SINR chung có xấp xỉ:
\begin{equation}
\label{eq:xap-xi-chung}
\gamma_c \approx \frac{p_c}{1 - p_c},
\end{equation}
ở đây đã chuẩn hóa $P_{\max} = 1$.

Trong đó: $p_c$ là tỷ lệ công suất luồng chung và $1 - p_c$ là tổng công suất luồng
riêng. Công thức cho thấy tăng $p_c$ làm việc giải mã luồng chung dễ hơn rất nhanh khi
$p_c$ tiến gần một, tạo thêm biên an toàn cho QoS của các UE chủ yếu dựa vào luồng chung.

Khi đó:
\begin{equation}
\label{eq:toc-do-chung-xap-xi}
R_c \approx -\log_2(1 - p_c).
\end{equation}

Trong đó: $R_c$ là tốc độ luồng chung trong xấp xỉ khi tạp âm nhỏ so với phần công suất
riêng tại UE nút cổ chai. Phần tốc độ chung tăng khi phần công suất riêng còn lại giảm;
đây là một nửa của cơ chế bù trừ giữa tốc độ chung và công suất luồng riêng.

Với UE $k^\star$ nhận luồng riêng chính:
\begin{equation}
\label{eq:toc-do-rieng-xap-xi}
R_{p,k^\star} \approx \log_2\!\left(\frac{g_{k^\star}(1 - p_c)}{\sigma^2}\right).
\end{equation}

Trong đó: $g_{k^\star}$ là độ lợi kênh của UE được ưu tiên và $1 - p_c$ là công suất
riêng còn lại. Khi tăng $p_c$, tốc độ riêng của UE mạnh giảm vì ngân sách riêng nhỏ đi;
tuy nhiên tốc độ chung tăng theo phần bù tương ứng.

Cộng hai thành phần cho trực giác:
\begin{equation}
\label{eq:tong-hai-thanh-phan}
R_c + R_{p,k^\star} \approx \log_2\!\left(\frac{g_{k^\star}}{\sigma^2}\right).
\end{equation}

Trong đó: hai hạng $\log_2(1 - p_c)$ triệt tiêu trong xấp xỉ nên kết quả còn phụ thuộc
chủ yếu vào độ lợi kênh của UE được chọn và tạp âm.

Công thức~\eqref{eq:tong-hai-thanh-phan} giải thích vì sao trong một vùng công suất chung
tương đối rộng, tăng $p_c$ có thể cải thiện khả năng cấp QoS cho UE yếu mà không làm tổng
tốc độ giảm nhiều. Đây là trực giác quan trọng để chọn một miền $p_c$ ưu tiên cao trong
bộ giải mã. Tuy nhiên, xấp xỉ chỉ đúng khi các giả thiết tỷ số tín hiệu trên tạp âm cao
và một luồng riêng chiếm ưu thế được thỏa.

\subsection{Vai trò của căn pha}
\label{ssec:vai-tro-can-pha}

Khi công suất luồng riêng gần tập trung vào một UE, việc tăng $g_{k^\star} =
|h^{\mathrm{eff}}_{k^\star}|^2$ trở nên rất có giá trị. Nếu bỏ qua đường trực tiếp và các
ràng buộc nhiều UE trong một bước trực giác, pha của từng phần tử nên được chọn sao cho
các tích $h^{*}_{r,k^\star,n} g_n e^{j\theta_n}$ có pha gần nhau. Khi đó, các thành phần
qua STAR--RIS cộng tăng cường thay vì triệt tiêu.

Hình~\ref{fig:ch2-can-pha} minh họa hiệu ứng này. Ở khung bên trái, các thành phần ghép
tầng có pha lệch nhau nên khi cộng tiếp nối, chúng tạo thành một đường gấp khúc gần khép
kín và véc-tơ tổng rất ngắn. Ở khung bên phải, sau khi mỗi phần tử được xoay về pha tham
chiếu, các thành phần nằm thẳng hàng và véc-tơ tổng dài hơn hẳn. Độ dài véc-tơ tổng chính
là $|h^{\mathrm{eff}}_{k^\star}|$ xuất hiện trong~\eqref{eq:tap-trung-cong-suat}.

\begin{figure}[htbp]
\centering
\resizebox{0.9\textwidth}{!}{\begin{tikzpicture}[font=\footnotesize,>=Stealth,scale=1.0]

% ---------- panel trai: pha lech ----------
\begin{scope}[shift={(0,0)}]
  \draw[gray!45,->] (-0.6,0)--(3.6,0);
  \draw[gray!45,->] (0,-1.9)--(0,1.9);
  \draw[blue!65,line width=1pt,->] (0,0) -- ++(20:0.75);
  \draw[blue!65,line width=1pt,->] (20:0.75) -- ++(100:0.75);
  \draw[blue!65,line width=1pt,->] ($(20:0.75)+(100:0.75)$) -- ++(175:0.75);
  \draw[blue!65,line width=1pt,->] ($(20:0.75)+(100:0.75)+(175:0.75)$) -- ++(250:0.75);
  \draw[blue!65,line width=1pt,->] ($(20:0.75)+(100:0.75)+(175:0.75)+(250:0.75)$) -- ++(320:0.75);
  \coordinate (T1) at ($(20:0.75)+(100:0.75)+(175:0.75)+(250:0.75)+(320:0.75)$);
  \draw[red!75!black,line width=1.5pt,->] (0,0) -- (T1);
  \node[font=\scriptsize,red!75!black,anchor=west] at (0.30,-0.30) {$\big|\sum_n\big|$ nhỏ};
  \node[font=\scriptsize,anchor=north] at (1.3,-2.15) {(a) Pha chưa căn chỉnh};
\end{scope}

% ---------- panel phai: da can pha ----------
\begin{scope}[shift={(6.6,0)}]
  \draw[gray!45,->] (-0.6,0)--(4.6,0);
  \draw[gray!45,->] (0,-1.9)--(0,1.9);
  \foreach \i in {0,1,2,3,4}{
     \draw[blue!65,line width=1pt,->] (\i*0.75,0.30) -- (\i*0.75+0.75,0.30);
  }
  \draw[red!75!black,line width=1.5pt,->] (0,-0.30) -- (3.75,-0.30);
  \node[font=\scriptsize,red!75!black,anchor=north] at (1.9,-0.38) {$\big|\sum_n\big|$ lớn};
  \node[font=\scriptsize,anchor=north] at (1.9,-2.15) {(b) Sau khi căn pha $\bar\theta_n=-\angle z_n$};
\end{scope}

\node[font=\scriptsize,align=center,anchor=north] at (4.2,-2.75)
  {Mỗi mũi tên là một thành phần ghép tầng $h^{*}_{r,k,n}g_n e^{j\theta_n}$ qua một phần tử STAR--RIS};
\end{tikzpicture}
}
\caption{Ý nghĩa của căn pha đối với độ lợi kênh hiệu dụng}
\label{fig:ch2-can-pha}
\figuresource{Tác giả xây dựng dựa trên \cite{Wu2019IRSBeamforming, DiRenzo2020SmartRadio} (2026)}
\end{figure}

Tuy nhiên, tối ưu chỉ theo UE mạnh nhất có thể làm giảm kênh của UE khác và làm $R_c$
giảm. Vì vậy, pha tham chiếu trong phương pháp đề xuất không cố định cho một UE theo quy
tắc cứng mà được tính từ hỗn hợp trọng số người dùng do Actor tạo ra. Sau đó, Actor còn
học phần hiệu chỉnh pha quanh tham chiếu để cân bằng mục tiêu RSMA và QoS.

\subsection{Hàm ý đối với thiết kế biểu diễn hành động}
\label{ssec:ham-y-thiet-ke}

Phân tích Chương~\ref{chap:mo-hinh-he-thong} dẫn đến ba hàm ý. Thứ nhất, không gian công
suất riêng cần dễ biểu diễn nghiệm tập trung cao; softmax có độ sắc lớn phù hợp hơn việc
buộc Actor tự tiến sát đỉnh đơn hình bằng đầu ra bão hòa. Thứ hai, tỷ lệ công suất luồng
chung nên được ưu tiên trong một vùng đủ cao để hỗ trợ QoS, nhưng không được tuyên bố là
ràng buộc vật lý của RSMA. Thứ ba, pha nên có một điểm xuất phát dựa trên căn chỉnh kênh
thay vì yêu cầu mạng học toàn bộ $[-\pi, \pi]$ từ đầu.

Ba hàm ý này tạo nên bộ giải mã có cấu trúc ở Chương~\ref{chap:phuong-phap}. Về mặt học
máy, đây là cách đưa thiên kiến quy nạp (inductive bias) vào chính không gian hành động. Về mặt vật lý, nó
tương đương với việc cung cấp cho mạng một bản đồ gần đúng về vùng nào đáng khám phá,
trong khi vẫn để mạng điều chỉnh theo CSI cụ thể.

\section{Giới hạn của mô hình}
\label{sec:gioi-han-mo-hinh}

Mô hình có một số giới hạn cần nêu rõ. Kênh Gaussian phức với hệ số tỷ lệ cố định không
mô tả đầy đủ suy hao theo khoảng cách, che khuất hay tương quan không gian. CSI hoàn hảo
bỏ qua sai số ước lượng. SIC lý tưởng bỏ qua lan truyền lỗi. Pha truyền và phản xạ được
điều khiển độc lập, trong khi một số kiến trúc phần cứng có ràng buộc ghép pha
\cite{Liu2021STARCoupledPhase, Wang2022STAROptimization}. Số UE cố định bằng bốn nên kết
quả chưa chứng minh khả năng tổng quát theo số người dùng. Những giới hạn này được xem là
phạm vi của nghiên cứu, không phải lỗi của thuật toán.

\section{Kết luận chương}
\label{sec:ch2-ket-luan}

Chương~\ref{chap:mo-hinh-he-thong} đã xây dựng đầy đủ mô hình vật lý và mô hình RSMA nhất
quán với mã nguồn. Phân tích cho thấy trong hệ SISO, nhiễu giữa nhiều luồng riêng có thể
làm SINR trở nên phụ thuộc chủ yếu vào tỷ lệ công suất thay vì độ lợi kênh. Luồng chung
cung cấp một cơ chế linh hoạt để duy trì QoS, còn căn pha STAR--RIS có thể tăng độ lợi
của UE được ưu tiên. Các kết quả này là cơ sở trực tiếp cho thiết kế bộ giải mã hành động
có cấu trúc trong Chương~\ref{chap:phuong-phap}.
